%%%%%%%%%%%%%%%%%%%%%%%%%%%%%%%%%%%%%%%%%
% Short Sectioned Assignment LaTeX Template Version 1.0 (5/5/12)
% This template has been downloaded from: http://www.LaTeXTemplates.com
% Original author:  Frits Wenneker (http://www.howtotex.com)
% License: CC BY-NC-SA 3.0 (http://creativecommons.org/licenses/by-nc-sa/3.0/)
%%%%%%%%%%%%%%%%%%%%%%%%%%%%%%%%%%%%%%%%%

% \documentclass[paper=a4, fontsize=11pt]{scrartcl} % A4 paper and 11pt font size
\documentclass[11pt, a4paper]{book}
\usepackage[T1]{fontenc} % Use 8-bit encoding that has 256 glyphs
\usepackage[utf8]{inputenc}
\usepackage{fourier} % Use the Adobe Utopia font for the document - comment this line to return to the LaTeX default
\usepackage{listings} % para insertar código con formato similar al editor
\usepackage[spanish, english, es-tabla]{babel} % Selecciona el español para palabras introducidas automáticamente, p.ej. "septiembre" en la fecha y especifica que se use la palabra Tabla en vez de Cuadro
\usepackage{url} % ,href} %para incluir URLs e hipervínculos dentro del texto (aunque hay que instalar href)
\usepackage{graphics,graphicx, float} %para incluir imágenes y colocarlas
\usepackage[gen]{eurosym} %para incluir el símbolo del euro
\usepackage{cite} %para incluir citas del archivo <nombre>.bib
\usepackage{enumerate}
\usepackage{hyperref}
\usepackage{graphicx}
\usepackage{tabularx}
\usepackage{booktabs}
\usepackage{threeparttable}

\usepackage[table,xcdraw]{xcolor}
\hypersetup{
	colorlinks=true,	% false: boxed links; true: colored links
	linkcolor=black,	% color of internal links
	urlcolor=cyan		% color of external links
}
\renewcommand{\familydefault}{\sfdefault}
\usepackage{fancyhdr} % Custom headers and footers
\pagestyle{fancyplain} % Makes all pages in the document conform to the custom headers and footers
\fancyhead[L]{} % Empty left header
\fancyhead[C]{} % Empty center header
\fancyhead[R]{Francisco Javier Bolívar Expósito} % My name
\fancyfoot[L]{} % Empty left footer
\fancyfoot[C]{} % Empty center footer
\fancyfoot[R]{\thepage} % Page numbering for right footer
%\renewcommand{\headrulewidth}{0pt} % Remove header underlines
\renewcommand{\footrulewidth}{0pt} % Remove footer underlines
\setlength{\headheight}{13.6pt} % Customize the height of the header

\usepackage{titlesec, blindtext, color}
\definecolor{gray75}{gray}{0.75}
\newcommand{\hsp}{\hspace{20pt}}
\titleformat{\chapter}[hang]{\Huge\bfseries}{\thechapter\hsp\textcolor{gray75}{|}\hsp}{0pt}{\Huge\bfseries}
\setcounter{secnumdepth}{4}
\usepackage[Lenny]{fncychap}


\begin{document}

	% Plantilla portada UGR
	\begin{titlepage}
\newlength{\centeroffset}
\setlength{\centeroffset}{-0.5\oddsidemargin}
\addtolength{\centeroffset}{0.5\evensidemargin}
\thispagestyle{empty}

\noindent\hspace*{\centeroffset}\begin{minipage}{\textwidth}

\centering
\includegraphics[width=0.9\textwidth]{logos/logo_ugr.jpg}\\[1.4cm]

\textsc{ \Large TRABAJO FIN DE GRADO}
\vspace{0.2cm}

\textsc{ GRADO EN INGENIERÍA INFORMÁTICA}
\vspace{1cm}

{\Huge\bfseries UltraStar: Song to txt \\}
\noindent\rule[-1ex]{\textwidth}{3pt}

\vspace{3.5ex}
{\large\bfseries Creación automática de archivos para Karaoke}
\end{minipage}

\vspace{2.5cm}
\noindent\hspace*{\centeroffset}
\begin{minipage}{\textwidth}
\centering

\textbf{Autor} 

{Francisco Javier Bolívar Expósito}
\vspace{2.5ex}

\textbf{Director} 

{Juan Julián Merelo Guervós}
\vspace{2cm}

\includegraphics[width=0.3\textwidth]{logos/etsiit_logo.png}\\[0.1cm]
\textsc{Escuela Técnica Superior de Ingenierías Informática y de Telecomunicación}\\
\textsc{---}\\
Granada, junio de 2022
\end{minipage}
\end{titlepage}


	% Plantilla prefacio UGR
	\thispagestyle{empty}

\begin{center}
\large\bfseries
UltraStar: \texttt{\normalfont{\large\bfseries Song2txt}}

Adaptación automática de canciones para juegos de karaoke
\end{center}

\begin{center}
Francisco Javier Bolívar Expósito
\end{center}


\vspace{0.5cm}
\noindent\textbf{Palabras clave}: \textit{Software libre, UltraStar, Python, Desarrollo ágil}
\vspace{0.7cm}

\noindent\textbf{Resumen}

La adaptación de canciones para videojuegos de karaoke es un proceso complejo y manual que requiere conocimientos musicales, lo que provoca que el número de canciones disponibles sea reducido. Aprovechando los avances en inteligencia artificial, gracias a los cuales es posible resolver problemas con sistemas informáticos que anteriormente requerían la intervención de un humano; este trabajo tiene el propósito de facilitar la adaptación de canciones avanzando en su automatización.

Para ello, utilizando un desarrollo ágil, se ha implementado una aplicación en Python que a partir del audio de una canción estima los tonos de cada nota interpretada por el cantante y los guarda en un archivo con el formato estándar UltraStar para videojuegos de karaoke, con una precisión superior a soluciones existentes. Adicionalmente, la librería desarrollada para trabajar con el formato UltraStar puede ser usada independientemente por otros proyectos.


\cleardoublepage

\begin{otherlanguage}{english}


\begin{center}
\large\bfseries 
\texttt{\normalfont\large\bfseries UltraStar: Song2txt}
	
Automatic adaptation of songs for karaoke games
\end{center}

\begin{center}
	\texttt{\normalfont Francisco Javier Bolívar Expósito}
\end{center}


\vspace{0.5cm}
\noindent\textbf{Keywords}: \textit{FLOSS, UltraStar, Python, Agile development}
\vspace{0.7cm}

\noindent\textbf{Abstract}

Adapting songs for karaoke video games is a complex and manual process that requires musical knowledge, which results in a limited number of songs available. Taking advantage of advances in artificial intelligence, thanks to which it is possible to solve problems with computer systems that previously required the intervention of a human; this work aims to facilitate the adaptation of songs by advancing in its automation.

To this end, using agile development, a Python application has been implemented that, from the audio of a song, estimates the tones of each note performed by the singer and saves them in a file with the standard \texttt{UltraStar} format for karaoke video games, with greater precision than existing solutions. Additionally, the library developed to work with the \texttt{UltraStar} format can be used independently by other projects.

\end{otherlanguage}

\cleardoublepage

\thispagestyle{empty}

\noindent\rule[-1ex]{\textwidth}{2pt}\\[4.5ex]

D. \textbf{Juan Julián Merelo Guervós}, Profesor del departamento de Arquitectura y Tecnología de Computadores.

\vspace{0.5cm}

\textbf{Informo:}

\vspace{0.5cm}

Que el presente trabajo, titulado \textit{\textbf{UltraStar: Song to txt}},
ha sido realizado bajo mi supervisión por \textbf{Francisco Javier Bolívar Expósito}, y autorizo la defensa de dicho trabajo ante el tribunal
que corresponda.

\vspace{0.5cm}

Y para que conste, expiden y firman el presente informe en Granada a junio de 2022.

\vspace{1cm}

\textbf{El/la director(a)/es: }

\vspace{5cm}

\noindent \textbf{Juan Julián Merelo Guervós}

\chapter*{Agradecimientos}

Aprovechando el fin de esta etapa académica, me gustaría agradecer a las personas que me han apoyado y gracias a las cuales soy quien soy.

En primer lugar, quiero agradecer el cariño de mi familia, especialmente a mis padres por su apoyo incondicional y a mi hermana por entenderme como nadie y por las pequeñas llamadas de 5 minutos que acaban durando horas.

También quiero dar las gracias a tantos profesores y compañeros de la universidad que no han dudado en prestar su ayuda, enseñándome tanto y compartiendo su pasión por este campo del conocimiento. A su vez, a la comunidad de software libre y a todas las personas que aportan desinteresadamente su granito de arena, confiando en llegar más lejos cooperando que compitiendo. 

Por último, a los \texttt{\normalfont\textit{Mákinas}} y los\texttt{\normalfont\textit{ Checolegas}}, la gente con la que he pasado algunos de los mejores momentos de esta etapa, con un gran corazón y con la que me siento en mi sitio.

Es imposible expresar todo lo que me habéis dado estos años y lo feliz que soy de teneros en mi vida. Un abrazo muy grande para todos.

\vspace*{1cm}



	% Índice de contenidos
	\tableofcontents

	% Índice de imágenes y tablas
	\listoffigures

	% Si hay suficientes se incluirá dicho índice
	\listoftables

	% Introducción 
	\chapter{Introducción}

Este proyecto es software libre, liberado con la licencia AGPLv3 \cite{agpl} y publicado en un repositorio público\footnote{\url{https://github.com/dipzza/ultrastar-song2txt}}.

\section{Motivación}
\label{sec:motivation}

El karaoke es un tipo de entretenimiento consistente en cantar junto a música grabada, normalmente con la ayuda de una pantalla en la que se muestra la letra de la canción e indicaciones de cuando hay que cantar cada parte.

Mucha gente disfruta del karaoke como forma de ocio o diversión, cumpliendo el sueño de ser cantante por unos momentos. Por esta razón, existen multitud de programas dirigidos a los entusiastas del karaoke.

Un tipo de estos programas son los videojuegos de karaoke, que ofrecen todavía más ayuda al mostrar el tono en el que se está cantando y como de lejos se encuentra este del correcto. Además, al medir la precisión con la que se ha cantado la canción permiten competir con otras personas o disfrutar buscando la mejora personal.

Sin embargo, los aficionados de los videojuegos de karaoke tienen que lidiar frecuentemente con un problema cuando quieren cantar su canción favorita: que no esté disponible.

Aunque el desarrollo de programas de software libre y gratuitos tanto de videojuegos de karaoke (\href{https://github.com/UltraStar-Deluxe/USDX}{\textit{UltraStar Deluxe}}, \href{https://github.com/Vocaluxe/Vocaluxe}{\textit{Vocaluxe}}, \href{https://github.com/performous/performous}{\textit{Performous}}...) como de herramientas para la creación de canciones (\href{https://github.com/sarutasan72/Yass}{\textit{Yass}}, \href{https://github.com/UltraStar-Deluxe/UltraStar-Creator}{\textit{UltraStar Creator}}...) ha ayudado al desarrollo de comunidades donde se crean y comparten canciones para videojuegos de karaoke, la mayoría de canciones siguen sin estar disponibles, especialmente aquellas más recientes o menos famosas.

Una causa de que la selección de canciones disponibles sea tan limitada es que el proceso de creación de canciones para videojuegos de karaoke es complejo, largo y requiere conocimientos musicales.


\section{Objetivo}

Este trabajo persigue ayudar en la creación de canciones para videojuegos de karaoke, reduciendo la barrera de entrada para personas sin conocimientos musicales y el tiempo que necesitan invertir aquellas con conocimientos musicales.


Para ello se pretende \textbf{mejorar o integrar los procesos de automatización existentes para la creación de canciones para videojuegos de karaoke.}

\section{Terminología}

\begin{itemize}
	\item{Nota. Representación musical de un sonido, con un tono constante y una duración determinada.}
	\item{Tono. Cualidad de los sonidos que permite su ordenación en una escala relacionada con la frecuencia.}
	\item{Frecuencia fundamental. Aquella que representa el tono simple más bajo presente en una nota.}
	\item{MIDI. Estándar técnico en el que se definen 128 notas por su número, con equivalencias a la nomenclatura estándar inglesa.}
	\item{Pulsación. Unidad básica en música para medir el tiempo. Una sucesión constante divide el tiempo en partes iguales.}
	\item{Pulsaciones por minuto. Unidad empleada para medir el ritmo en música. Equivale al número de pulsaciones que caben en un minuto.}
	\item{Fonema. Unidad sonora mínima que puede distinguir una palabra de otra en un lenguaje dado.}
	\item{Software libre. Aquel cuyo código fuente puede ser estudiado, modificado, utilizado libremente para cualquier propósito y redistribuido con cambios o mejoras.}
	\item{Software privativo. Aquel del cual no existe una forma libre de acceso a su código fuente, en el que uno o varios de los derechos de usar, modificar, compartir modificaciones o el software están restringidos.}
\end{itemize}

	% Descripción del problema y hasta donde se llega
	\chapter{Descripción del problema}

Con el fin de facilitar el proceso de añadir canciones a videojuegos de karaoke, se pretende avanzar en la automatización de las tareas relacionadas.

Independientemente del formato utilizado para guardar las canciones, las principales tareas consisten en:

\begin{itemize}
	\item{Encontrar e introducir metadatos, como el título, el artista, las pulsaciones por minuto, etc.}
	\item{Obtener la letra de la canción.}
	\item{Definir las notas interpretadas por el cantante original, con la siguiente información asociada a cada una:}
	\begin{itemize}
		\item{El tono del sonido.}
		\item{El intervalo de tiempo en el que sucede.}
		\item{Los fonemas de la letra correspondientes.}
	\end{itemize}
\end{itemize}

Algunas de las tareas son más complejas o tardan más en completarse. Definir la información de las notas en general es un proceso tedioso debido a la gran cantidad de notas que pueden ser cantadas en una canción (por ejemplo, la canción \texttt{'Rolling in the Deep' de Adele} requiere la definición de más de 500 notas individuales) y, en particular, la asignación de los tonos es especialmente compleja porque requiere de oído musical y un proceso de prueba y error.

¿Debe darse prioridad a la automatización de estas tareas, por ser las más complejas? Es un argumento a tener en cuenta, sin embargo, es importante valorar las soluciones ya existentes para decidir el enfoque del trabajo, analizadas en el capítulo \ref{ch:artstate} sobre el estado del arte.

	% Estado del arte
	% 	1. Crítica al estado del arte
	% 	2. Propuesta
	\chapter{Estado del arte}

El software libre y sus licencias \cite{gplv3} ha permitido llevar a cabo una expansión del
aprendizaje de la informática sin precedentes.

	
	\chapter{Planificación}

\section{Metodología utilizada}

\section{Análisis de usuarios}

Para empatizar con los usuarios que acabarán usando el proyecto es necesario definirlos claramente, pudiendo así enfocar el proyecto a resolver sus problemas.

Con esta intención, utilizando la metodología de personas basadas en roles \cite{role-personas}, se han definido los siguientes roles que describen a un grupo de personas que realizan la misma función.

\begin{enumerate}	
	\item{Miembro del tribunal}
	\begin{itemize}
		\item{Misión: Evaluar los actos de defensa y los informes presentados en los trabajos de fin de grado, teniendo en cuenta si estas muestran que el estudiante ha desarrollado las técnicas y habilidades propias para obtener el título de grado.}
		\item{Nivel de conocimiento: No se pueden hacer suposiciones sobre las habilidades de un miembro del tribunal, según su rama de conocimiento pueden tener un conocimiento altamente avanzado sobre los temas tratados o uno muy básico.}
		\item{Contexto en el que realiza la función: La función es parte de su trabajo. Para la evaluación del acto de defensa prestan normalmente toda su atención, al leer la memoria también suele haber cierto nivel de concentración, aunque este puede ser menor al poder tener sobrecarga de obligaciones y estrés.}
		\item{Productos de los que es usuario: Informe y el acto de defensa.}
	\end{itemize}
	
	\item{Creador de canciones para videojuegos de karaoke}
	\begin{itemize}
		\item{Misión: Adaptar canciones para poder utilizarlas en los videojuegos de karaoke.}
		\item{Nivel de conocimiento: Manejo medio de ordenadores, el sistema operativo que usan en estos puede ser \textit{Windows}, \textit{Linux} o \textit{macOS}. Importante diferenciar entre aquellos que tienen conocimientos musicales y los que no.}
		\item{Contexto en el que realiza la función: Empleando su tiempo libre, por lo que pueden dejarla a medias o tener un tiempo limitado para hacerla.}
		\item{Productos de los que es usuario: Las aplicaciones desarrolladas.}
	\end{itemize}
	
		\item{Programador}
	\begin{itemize}
		\item{Misión: Desarrollar con el código del proyecto, ya sea contribuyendo a este o reutilizando partes en otros proyectos.}
		\item{Nivel de conocimiento: Conocimiento experto sobre herramientas de desarrollo y colaboración.}
		\item{Contexto en el que realiza la función: Puede ser un su tiempo libre o como parte de sus tareas, pero suele ser con cierta dedicación sin hacer otras tareas a la vez.}
		\item{Producto de los que es usuario: La organización del proyecto, el código de este y la documentación.}
	\end{itemize}
\end{enumerate}




\section{Temporización}

\section{Seguimiento del desarrollo}


	% Análisis del problema
	% 1. Análisis de requisitos
	% 2. Análisis de las soluciones
	% 3. Solucion propuesta
	% 4. Análisis de seguridad
	\chapter{Análisis del problema}
 


	% Desarrollo bajo sprints: 
	% 	1. Permitir registros y login de usuarios
	% 	2. Desarrollo del sistema de incidencias
	% 	3. Desarrollo del sistema de denuncias administrativas y accidentes
	% 	4. Desarrollo del sistema de croquis
	%   5. Instalación de la aplicación de manera automática
	\chapter{Implementación}

La implementación del software se ha dividido en hitos. Estos, han sido definidos en Github
y cada uno de ellos contiene un grupo de \textit{issues} que se corresponden con las distintas
mejoras que se han ido incorporando al software a lo largo de su desarrollo.\\



	% Presupuesto

	% Conclusiones
	\chapter{Conclusiones y trabajos futuros}

\section{Conclusiones}

Con el trabajo realizado se ha conseguido avanzar en la creación automática de canciones para el formato UltraStar, cumpliendo el objetivo fijado para el proyecto.

Se ha contribuido al estado del arte, desarrollando la primera biblioteca pública en Python para trabajar con el formato UltraStar e implementando una aplicación para la asignación automática de tonos, que respecto a las existentes previamente: 

\begin{itemize}
	\item{Usa técnicas para la estimación de la frecuencia fundamental de mayor precisión, gracias a los avances recientes en redes neuronales.}
	\item{Provee una instalación más accesible con un empaquetado moderno, compatible con los principales sistemas operativos de escritorio y sin necesidad de instalar manualmente programas de terceros.}
	\item{Aporta documentación sobre los métodos empleados y las decisiones tomadas, tanto en esta memoria como en el propio repositorio donde se aloja la aplicación.}
\end{itemize}


Personalmente, la realización del trabajo ha servido no solo para demostrar habilidades adquiridas durante el grado, sino también como un vehículo para seguir aprendiendo sobre las apasionantes tareas propias de la creación de software eligiendo las soluciones más apropiadas.

Por último, me gustaría destacar la importancia del software libre, sin el cual este trabajo no habría sido posible, ya que ha sido crucial tanto en el aprendizaje como en el desarrollo, y, por tanto, la importancia de liberar el propio trabajo para facilitar futuras aportaciones.

\section{Trabajos futuros}

Alcanzar la automatización completa de las tareas descritas en la sección \ref{sec:problemdescription} está todavía muy lejos. Por ello, en contribuciones futuras se insta a seguir las siguientes líneas de trabajo:

\begin{itemize}
	\item{Desarrollar herramientas para obtener los metadatos de canciones a partir de servicios públicos en línea y guardarlos en el formato UltraStar (posiblemente haciendo uso de la librería implementada en este trabajo).}
	\item{Investigar e integrar soluciones para mejorar la precisión o la velocidad en la estimación de tonos, como:
	\begin{itemize}
		\item{Reducir la influencia de predicciones atípicas, ya sea haciendo la mediana de las predicciones en vez de la media u obviándolas si su desviación típica es demasiado alta.}
		\item{Utilizar técnicas del estado del arte para separar la voz del cantante del resto del audio previamente a la predicción, reduciendo así el ruido.}
		\item{Incrementar la precisión con nuevos modelos de redes neuronales o técnicas posteriores al desarrollo del proyecto, como\texttt{ HarmoF0}\cite{harmof0}.}
	\end{itemize}
	}
\end{itemize}


	% Trabajos futuros


	
	\bibliography{bibliografia}
	\bibliographystyle{plain}
	
\end{document}

